% Nouveau EGA V (cote 25): archive cote 25, site page 153 (PDF page 154).
% typescript. Draft compared: none.
% AI-transcribed by Claude Opus 5.5 (claude-opus-5-5, adaptive thinking, effort high), Message Batches API, batch ega25b (msgbatch_01A7n5YuXubUPNK1zsepi8RG), 2026-11-13; original AI transcription, subsequently checked against both supplied scan parts; see PAGES_CHECKED.tsv and CHECK_DECISIONS.tsv.
\scanpage{25/153}
\mspage{\hand{14}}

notions, mis à part \var{\xout{\unsure{une}}}{}{typewriter strike-out}{ega25_25_153-v0001}\var{\inserted[typed]{la}}{la}{insertion}{ega25_25_153-v0002} condition dimensionelle = la condition de
$\underline{O}_X$-régularité, le caractère de x se voit \var{\xout{\unsure{sur}}}{}{typewriter strike-out}{ega25_25_153-v0003} e\unsure{b|h} fait sur la nature du
point x dans V(u), i.e. sur l'anneau local $\underline{O}_{V(u),x}$ (éventuellement consi-%
déré comme algèbre sur k). Lorsque (dans le cas d'un corps de base k)
x est rationnel sur k, la notion relative à k coincide avec la notion
absolue, comme on voit \var{\inserted{en \unsure{procédant|précédemment}}}{en \unsure{procédant|précédemment}}{insertion}{ega25_25_153-v0004} comme dans le cas où $\underline{E}$ est inversible. Enfin,
lorsque X est lisse sur \var{\over{S}{Y}}{Y}{writing over}{ega25_25_153-v0005}, et u: $\underline{E}\to\underline{O}_X$ est encore un homomorphisme avec
$\underline{E}$ loc libre de type fini, on définit les notions correspondantes relative-%
ment à S, en regardant fibre par fibre. On remarquera que si x est
non supersingulier relativement à \var{\over{S}{Y}}{Y}{writing over}{ega25_25_153-v0006}, alors V(u) est plat sur S \var{\xout{relativ}}{}{typewriter strike-out}{ega25_25_153-v0007} en
x (car u est $\underline{O}_X$-régulier relativement à \var{\over{S}{Y}}{Y}{writing over}{ega25_25_153-v0008} ...).

\var{\margin{$\Vert$}}{\placedmargin{$\Vert$}}{marginal note placed}{ega25_25_153-v0009} Bien entendu (à signaler tout au moins en r emarque après les
définitions en forme) on pouvait tout aussi bien partir d'une section u
d'un Module loc libre $\underline{E}$ , les définitions précédentes s'appliquent en
interprétant u comme un homomorphisme $\overset{\hand{\scriptscriptstyle\vee}}{\underline{E}}\hand{\longrightarrow}\underline{O}_X$ .

Lorsque X est lisse sur \var{\over{S}{Y}}{Y}{writing over}{ega25_25_153-v0010}, u:$\underline{E}\hand{\longrightarrow}\underline{O}_X$ est donné avec $\underline{E}$ loc libre
de rang m, alors l'ensemble des \var{\xout{points}}{}{typewriter strike-out}{ega25_25_153-v0011} \var{\inserted[typed]{zéros de u}}{zéros de u}{insertion}{ega25_25_153-v0012} singuliers rel à S est l'ensembl
sous-jacent à un sous-préschéma fermé $V(u)^{\mathrm{sing}} = V^1$ bien déterminé de
$V(u)=V^o$ , obtenu ainsi: on considère $d^1u$ section de $P^1_{X/\var{\over{S}{Y}}{Y}{writing over}{ega25_25_153-v0013}}(\overset{\hand{\scriptscriptstyle\vee}}{\underline{E}})$ , sa
restriction à $V^o$ peut alors s'interpréter comme un homomorphisme
\[ \underline{E}\otimes\underline{O}_{V^o} \hand{\longrightarrow} \underline{\hand{\Omega}}^1_{X/\var{\over{S}{Y}}{Y}{writing over}{ega25_25_153-v0014}}\otimes\underline{O}_{V^o} \]
et $V(u)^{\mathrm{sing}}=V^1$ désignera le préschéma des zéros de \var{\xout{\unsure{ce dernier. Rem}}}{}{typewriter strike-out}{ega25_25_153-v0015}
\var{\xout{Considérant de même $d^2$ section de $P^2_{X/S}(\overset{\hand{\scriptscriptstyle\vee}}{\underline{E}})$ puis $d^2$ $V^1$ on}}{}{typewriter strike-out}{ega25_25_153-v0016}
\var{\xout{trouve un homomorphisme canonique}}{}{typewriter strike-out}{ega25_25_153-v0017} l'homomorphisme correspondant
\[ \hand{\bigwedge}^m\underline{E}\otimes\underline{O}_{V^o} \hand{\to} \underline{\hand{\Omega}}^m_{X/\var{\over{S}{Y}}{Y}{writing over}{ega25_25_153-v0018}}\otimes\underline{O}_{V^o} \]
i.e. d'une section bien déterminée de $\underline{\hand{\Omega}}^m_{X/\var{\over{S}{Y}}{Y}{writing over}{ega25_25_153-v0019}}\otimes\mathrm{dét}(E)^{-1}\otimes\underline{O}_{V^o}$ .

\mspage{\hand{153}}
